%% =========================================================================
\section{Variants of the Conjecture}
\label{sec:variants}
%% =========================================================================

The Euclid--Mullin construction makes sense with any factor-selection
rule: at each step, $\Prod{2}(n)+1$ has a non-empty set of prime
factors, and the rule picks one.  Mullin's Conjecture is specific to
the $\minFac$ rule.  This section collects the known variants---some
easy, some provably incomplete, one almost closed---and uses them to
locate standard MC within a broader landscape.

\subsection{The Selection Rule Landscape}
\label{sec:selection-landscape}

Fix an initial prime $p_0 = 2$.  Given a \emph{selection rule}
$\sigma$ that, for any integer $N \geq 2$, chooses a prime factor
$\sigma(N) \mid N$, define:
\begin{align}
  P_\sigma(0) &= 2, &
  P_\sigma(n{+}1) &= P_\sigma(n) \cdot \sigma(P_\sigma(n)+1).
\end{align}
The selection rule determines the sequence; the running product
$P_\sigma(n)$ determines the sequence history.
Five rules are studied in the formalization:

\smallskip
\begin{center}
\begin{tabular}{@{}l@{\quad}l@{\quad}l@{}}
\toprule
\textbf{Selection rule} $\sigma$ & \textbf{MC status} & \textbf{Key property} \\
\midrule
Random (uniform) & Easy (mixing) & Independence \\
$\minFac$ (standard EM) & \textbf{Open} (DSL) & Deterministic, lethal \\
$\{\minFac, \operatorname{secondMinFac}\}$ & \textbf{Open} (UFDStrong) & Two-point diversity \\
Steered (Booker~\cite{BoSa2012}) & True (by construction) & Adversarial \\
$\mathrm{maxFac}$ (Cox--vdP~\cite{CoxVdP1968}) & False & Misses primes \\
\bottomrule
\end{tabular}
\end{center}

\smallskip\noindent
The landscape reveals a tension: the random rule is the easiest to
analyze (mixing theory) but the worst at selection (probability
$1/|F|$ per hit); $\minFac$ is the best at selection (lethal hits)
but the hardest to analyze (deterministic, no independence); and the
$\{\minFac, \operatorname{secondMinFac}\}$ variant sits between them,
offering a two-point spectral contraction that is provably sufficient
under a single open hypothesis.

\subsection{The Random-Factor Variant}
\label{sec:random-variant}

Consider a variant where, instead of the smallest prime factor, one
picks a \emph{random} prime factor of~$\Prod{2}(n)+1$ at each step.
MC follows from classical mixing theory.  For any target prime~$p$,
the residue $r_n = \Prod{2}(n) \bmod p$ performs a multiplicative walk
on $(\ZZ/p\ZZ)^\times$.  Each multiplier is a random element of the
group; once the multipliers generate the full group (which PRE
guarantees, Theorem~\ref{thm:pre}), the walk is a genuine random walk
with full support.  By Diaconis--Shahshahani equidistribution, the walk
mixes in $O(\log p)$ steps: $r_n = -1$ occurs with frequency
$\to 1/(p{-}1)$---infinitely often with probability~1.  When $p$
divides the Euclid number, the random rule selects~$p$ with probability
$1/|F(S_n)|$, but cofinal hitting provides infinitely many chances,
so~$p$ is eventually selected.

\paragraph{Comparison with $\minFac$.}
The actual EM sequence uses
$\seq{2}(n{+}1) = \minFac(\Prod{2}(n){+}1)$, a \emph{deterministic}
function of the walk position.  Compared to the random-factor variant,
$\minFac$ is strictly better in every measurable dimension:

\smallskip
\begin{center}
\begin{tabular}{l@{\quad}l@{\quad}l}
\toprule
 & \textbf{Random factor} & $\boldsymbol{\minFac}$ \\
\midrule
\textbf{Selection when $q \mid P(n){+}1$}
  & Probability $1/|F(S_n)|$ & Probability~1 \\
\textbf{Accumulator growth}
  & Faster (random $\geq \minFac$) & Slowest possible \\
\textbf{Hits on $-1$ needed}
  & Infinitely many & \textbf{One} (lethal hit) \\
\bottomrule
\end{tabular}
\end{center}

\smallskip\noindent
\emph{Selection efficiency.}\enspace Past the sieve gap, 0-accessibility
guarantees: if $q \mid \Prod{2}(n)+1$ then
$q = \minFac(\Prod{2}(n)+1)$---selection is certain, not
probabilistic.  \emph{Growth rate.}\enspace Since
$\minFac(N) \leq p$ for any factor~$p$ of~$N$, the accumulator
$\Prod{2}(n)$ grows minimally under the $\minFac$ rule, giving
exponentially more steps before it becomes astronomically large---more
chances for~$q$ to divide.  \emph{Lethality.}\enspace The inductive
bootstrap shows one single hit suffices: $\MC({<}\,q)$ ensures~$q$
is minimal whenever it divides, so hitting $-1$ once means~$q$
enters the sequence.  The random variant requires infinitely many
hits to overcome low per-hit selection probability.

\emph{MinFac is strictly better in every measurable dimension.
So why can't we prove it works?}

\paragraph{The independence trade-off.}
The answer is surgical: the random variant possesses a single property
that $\minFac$ lacks, and that property is the \emph{only thing the
proof uses}.

The property is \textbf{independence of multiplier from walk
position}.  In the random-factor variant, the multiplier at step~$n$
is selected independently of which residue class
$\Prod{2}(n) \bmod q$ occupies.  Walk position and multiplier are
independent random variables, and independence is the engine that
drives mixing.  Diaconis--Shahshahani's theorem applies whenever the
multipliers are i.i.d.\ from a distribution whose support generates
the group.

In the $\minFac$ variant, the multiplier is
$\minFac(\Prod{2}(n){+}1) \bmod q$---a deterministic function of the
full integer~$\Prod{2}(n)+1$.  The CRT multiplier invariance
(\lean{EM/Group/CRT.lean}{crt\_multiplier\_invariance}) shows
that the outcome is independent of $\Prod{2}(n) \bmod q$ at the
\emph{population} level: for generic integers, knowing the mod-$q$
residue does not predict $\minFac$.  But for the specific orbit~$n=2$,
the multiplier is a deterministic value, and mixing theorems for
deterministic sequences cannot be invoked.

The random variant thus sacrifices selection efficiency (it may
miss~$q$ even when $q$ divides) in exchange for statistical
independence.  The $\minFac$ variant has perfect selection but no
independence:

\smallskip
\begin{center}
\begin{tabular}{l@{\quad}l@{\quad}l}
\toprule
 & \textbf{Random factor} & $\boldsymbol{\minFac}$ \\
\midrule
Selection & Bad ($1/k$ chance) & Perfect \\
Mixing & Random walk theorem & Unproved \\
\bottomrule
\end{tabular}
\end{center}

\smallskip\noindent
The random variant's proof: mixing $\to$ cofinal hitting $\to$
infinitely many chances compensate for low per-hit probability.
The $\minFac$ variant needs only one hit---absurdly less than the
random variant achieves---but has no way to produce even that one
without mixing.

\paragraph{The stochastic domination failure.}
\label{sec:stoch-dom}
The natural instinct is to argue by \emph{stochastic domination}:
if the random walk visits~$-1$ cofinally, and $\minFac$ is ``better
in every way,'' then $\minFac$ should also visit~$-1$.  This argument
fails because the two processes \emph{diverge after one step}.

At step~$n$, both processes face the same Euclid number
$\Prod{2}(n)+1$ with factor set~$F(S_n)$.  The random variant picks a
uniformly random element of~$F(S_n)$; the $\minFac$ variant picks the
minimum.  After this step, the two processes have different
accumulators (they multiplied by different primes), so they face
different Euclid numbers at the next step.  The processes explore
completely different trajectories.  Domination fails because the
two walks live on different probability spaces after step~1; we
cannot compare them step-by-step.

A subtler comparison uses \emph{visit measures}.  Define
$\mu_{\mathrm{rand}}(a) = \lim_{N} \frac{1}{N}|\{n < N :
w_q^{\mathrm{rand}}(n) = a\}|$ and analogously
$\mu_{\min}(a)$ for the $\minFac$ walk.  For the random variant,
$\mu_{\mathrm{rand}} = $ uniform on $(\ZZ/q\ZZ)^\times$ by mixing.
For $\minFac$, CME/CCSB asserts $\mu_{\min} \approx$ uniform, and
DSL asserts this follows from PE\@.  The comparison
$\mu_{\min}(-1) \geq 1/(q{-}1)$ is precisely what needs proof, and is
the content of DSL\@.

\subsection{The $\{\minFac, \operatorname{secondMinFac}\}$ Variant}
\label{sec:minfac-secondminfac}

The two-point variant interpolates between the random and deterministic
extremes.  At each step, instead of a single deterministic choice or a
fully random one, two specific factors are available:
$\minFac(\Prod{2}(n)+1)$ and $\operatorname{secondMinFac}(\Prod{2}(n)+1)$.
A decision sequence $\sigma : \NN \to \mathrm{Bool}$ selects between
them; the all-true sequence recovers the standard EM walk.

\paragraph{The $\varepsilon$-walk.}
The formalization defines a
\textbf{self-consistent $\varepsilon$-walk}
(\lean{EM/Advanced/VanishingNoiseVariant.lean}{epsWalkProd})
in which the chosen factor---either $\minFac(P(n){+}1)$ or
$\operatorname{secondMinFac}(P(n){+}1)$---becomes the actual multiplier,
updating the accumulator and producing a new Euclid number at the next
step.
The function $\operatorname{secondMinFac}(n)$ (second-smallest prime
factor) is defined and proved prime, divisor of~$n$, and
$\geq \minFac(n)$
(\lean{EM/Advanced/VanishingNoiseVariant.lean}{secondMinFac\_prime},
\lean{EM/Advanced/VanishingNoiseVariant.lean}{secondMinFac\_dvd},
\lean{EM/Advanced/VanishingNoiseVariant.lean}{minFac\_le\_secondMinFac}).

This self-consistent walk generates a \emph{binary branching tree}
of depth~$N$ with $2^N$ leaves, where each internal node has its own
path-dependent factor set.  The \emph{tree character sum}
(\lean{EM/Advanced/VanishingNoise.lean}{treeCharSum}),
a weighted sum of $\chi(\text{leaf position})$ over all leaves, is proved
$\leq 1$ in norm.  Unlike the product multiset of the stochastic MC, this sum
does \textbf{not} factorize as $\prod_k \bar\chi(S_k)$: the left and
right subtrees start from different accumulators, so their character sums
differ.  The open hypothesis
\defname{TreeContractionHypothesis}
asserts that the tree character sum tends to zero for all nontrivial~$\chi$;
it is strictly weaker than DSL but strictly stronger than the product approximation.

\paragraph{Non-self-consistent variant MC.}
An independent route bypasses the tree contraction problem entirely.
At each EM step~$n$, form the two-element unit set
$\{\minFac(P(n){+}1),\, \operatorname{secondMinFac}(P(n){+}1)\} \bmod q$
in $(\ZZ/q\ZZ)^\times$.  When one factor equals~$q$
(a non-unit), fall back to $\operatorname{Finset.univ}$, whose
mean character value is~$0$ by orthogonality---giving perfect contraction
(\lean{EM/Advanced/VanishingNoiseVariant.lean}{meanCharValue\_univ\_eq\_zero}).
Under \defname{UFDStrong}---spectral gaps of the padded unit sets
are non-summable for every nontrivial~$\chi$---the
sparse product contraction machinery gives vanishing averaged
products, and the Fourier counting argument yields
path existence.  Concretely:

\begin{theorem}[\lean{EM/Advanced/VanishingNoiseVariant.lean}{variant\_mc\_from\_ufd\_strong\_proved}]
\label{thm:variant-mc}
For every prime $q \geq 3$, if $\mathrm{UFDStrong}(q)$ holds, then
every element of $(\ZZ/q\ZZ)^\times$ is representable as a product
of selections from $\{\minFac,\operatorname{secondMinFac}\}$ of
consecutive Euclid numbers.  In particular, some selection path
hits~$-1$.
\end{theorem}

The verified chain is:
\[
  \mathrm{UFDStrong}(q)
  \;\Longrightarrow\;
  \text{vanishing products}
  \;\Longrightarrow\;
  \text{path existence}
  \;\Longrightarrow\;
  \text{VariantMC}.
\]
Every arrow is machine-checked with zero \code{sorry}.

\paragraph{The gap: UFDStrong.}
\defname{UFDStrong}$(q)$ asserts: for each nontrivial character
$\chi$ on $(\ZZ/q\ZZ)^\times$, the spectral gaps
$1 - \|\bar\chi(\text{paddedUnitSet}_n)\|$ are non-summable.
The hypothesis splits into two cases:
\begin{enumerate}[nosep]
\item \emph{Fallback case}: if infinitely many steps have
  $\minFac \equiv \operatorname{secondMinFac} \pmod{q}$ (or one
  equals~$q$), the padded set falls back to the full group, giving
  perfect contraction (gap~$= 1$) at each such step.
  Non-summability is immediate.
  \textbf{Proved}:
  \lean{EM/Advanced/VanishingNoiseVariant.lean}{ufd\_fallback\_not\_summable}.
\item \emph{Diversity case}: if cofinitely many steps have two
  distinct units mod~$q$, one needs the character values at those
  units to differ infinitely often---i.e., the ratio
  $\minFac/\operatorname{secondMinFac}$ does not persistently lie in
  $\ker(\chi)$.  This is a genuine number-theoretic question about
  Euclid number factorizations.  \textbf{Open}:
  \code{UFDImpliesUFDStrong}.
\end{enumerate}

\paragraph{Variant MC does not imply standard MC.}
The selection~$\sigma$ that achieves hitting may differ per
prime~$q$, and the virtual accumulator
$\prod_k \text{factor}_{\sigma(k)}$ is not the EM accumulator
$P(n)$.  The gap is the same orbit-specificity barrier
(Dead End~\#130,
\lean{EM/FunctionField/CyclicWalkCoverage.lean}{alternating\_walk\_misses\_two}):
generation does not imply coverage for deterministic walks.

\paragraph{Three routes to UFDStrong.}
The formalization identifies a hierarchy of progressively weaker
sufficient conditions for UFDStrong$(q)$:

\begin{enumerate}[nosep]
\item \textbf{Quantitative escape} (\defname{MinFacRatioEscape}):
  for each nontrivial~$\chi$, there exists $\delta > 0$ such that
  infinitely many steps have spectral gap $\geq \delta$.
  Since a non-summable sequence with a positive lower bound infinitely
  often is non-summable, this implies UFDStrong directly
  (\lean{EM/Advanced/VanishingNoiseVariant.lean}{ratio\_escape\_implies\_ufdStrong}).

\item \textbf{Qualitative escape} (\defname{MinFacRatioEscapeQual}):
  for each nontrivial~$\chi$, infinitely many steps have
  $\operatorname{card}(\text{rawTwoPointUnitSet}) \geq 2$ \emph{and}
  $\chi(\minFac) \neq \chi(\operatorname{secondMinFac})$.
  This implies quantitative escape via a finite-group pigeonhole
  argument: on a finite group, the gap function takes finitely many
  values, so being positive infinitely often gives a uniform positive
  lower bound
  (\lean{EM/Advanced/VanishingNoiseVariant.lean}{qual\_implies\_quant}).

\item \textbf{Orbit-level minFac equidistribution} (\defname{OrbitMFRE}):
  the residues $\minFac(\Prod{2}(n){+}1) \bmod q$ are equidistributed
  among the units of $(\ZZ/q\ZZ)^\times$ in Cesàro density.
  This is the weakest condition in the hierarchy: equidistributed
  residues prevent persistent chi-degeneracy.  The bridge
  OrbitMFRE $\Rightarrow$ MinFacRatioEscapeQual is \textbf{open}
  (\code{OrbitMFREImpliesEscapeQual}).
\end{enumerate}

The assembled chain is:
\[
  \underbrace{\mathrm{OrbitMFRE}}_{\text{open}}
  \;\xRightarrow{\text{open bridge}}\;
  \mathrm{EscapeQual}
  \;\xRightarrow{\text{proved}}\;
  \mathrm{EscapeQuant}
  \;\xRightarrow{\text{proved}}\;
  \mathrm{UFDStrong}
  \;\xRightarrow{\text{proved}}\;
  \mathrm{VariantMC}.
\]
Every arrow except the first is machine-checked with zero \code{sorry}.
The most concrete attack target is MinFacRatioEscapeQual: one needs to
show that $\minFac$ and $\operatorname{secondMinFac}$ do not
persistently agree modulo every character.  OrbitMFRE faces the same
orbit-specificity barrier as DSL for standard MC.

\paragraph{Stochastic Two-Point MC.}
The non-self-consistent machinery yields a clean quantitative statement.
Define the \defname{StochasticTwoPointMC} problem: under UFDStrong$(q)$
with $q \geq 3$, does every element of $(\ZZ/q\ZZ)^\times$ appear with
positive count in the product multiset of padded unit sets?

\begin{theorem}[\lean{EM/Advanced/VanishingNoiseVariant.lean}{stochastic\_two\_point\_mc\_proved}]
\label{thm:stochastic-two-point}
$\mathrm{UFDStrong}(q) \;\Longrightarrow\; \mathrm{StochasticTwoPointMC}(q)$.
\end{theorem}

\noindent The product multiset grows exponentially
($\geq 2^N$ paths for $q \geq 3$,
\lean{EM/Advanced/VanishingNoiseVariant.lean}{productMultiset\_card\_ge\_two\_pow}).
A \defname{fairCoin} schedule ($\varepsilon = 1/2$ at every step) gives
the \defname{fairTreeCharSum}---the tree character sum at constant
$1/2$---which is bounded by~$1$ in norm
(\lean{EM/Advanced/VanishingNoiseVariant.lean}{fairTreeCharSum\_norm\_le\_one}).
TreeContractionHypothesis implies the weaker
\defname{TreeContractionAtHalf}
(\lean{EM/Advanced/VanishingNoiseVariant.lean}{treeContraction\_of\_hypothesis}),
and at $\varepsilon = 0$ the tree degenerates to the deterministic
$\minFac$ path
(\lean{EM/Advanced/VanishingNoiseVariant.lean}{treeCharSum\_at\_zero\_step}).

The gap between the non-self-consistent result (proved) and a
\emph{self-consistent} stochastic MC is the \textbf{selection coherence}
problem: the non-self-consistent product multiset uses factor sets from
the actual EM walk (which always picks $\minFac$), so the factor sets
are fixed and the product factorizes; in the self-consistent tree,
picking $\operatorname{secondMinFac}$ at step~$n$ changes the
accumulator, hence the factor set at step~$n{+}1$, preventing
factorization.

\paragraph{Faithful character escape.}
A character $\chi$ on $(\ZZ/q\ZZ)^\times$ is \defname{faithful} if it
is injective (equivalently, $\ker(\chi) = \{1\}$).  For faithful
characters, whenever the raw two-point unit set has two distinct
elements, their $\chi$-values are automatically distinct---giving a
positive spectral gap at every non-fallback step, with no
number-theoretic hypothesis.  Combined with the fallback case (gap $= 1$
at fallback steps, already proved), this yields unconditional
non-summability:

\begin{theorem}[\lean{EM/Advanced/VanishingNoiseVariant.lean}{faithful\_character\_escape}]
\label{thm:faithful-escape}
For every prime $q \geq 3$ and every faithful nontrivial character
$\chi$ on $(\ZZ/q\ZZ)^\times$:
\[
  \sum_n \bigl(1 - \|\bar\chi(\text{paddedUnitSet}_n)\|\bigr)
  = +\infty.
\]
Unconditional---zero open hypotheses.
\end{theorem}

\noindent For groups of prime order, every nontrivial character is
faithful (by Lagrange: the kernel is a subgroup of a prime-order group,
hence trivial or the whole group;
\lean{EM/Advanced/VanishingNoiseVariant.lean}{prime\_order\_nontrivial\_is\_faithful}).
Since $|(\ZZ/3\ZZ)^\times| = 2$ is prime, UFDStrong$(3)$ holds
\emph{unconditionally}
(\lean{EM/Advanced/VanishingNoiseVariant.lean}{ufdStrong\_three}), and
the full chain closes:

\begin{theorem}[\lean{EM/Advanced/VanishingNoiseVariant.lean}{variant\_mc\_three\_unconditional}]
\label{thm:variant-mc-3}
$\mathrm{StochasticTwoPointMC}(3)$ holds with zero open hypotheses:
every element of $(\ZZ/3\ZZ)^\times$ appears with positive count
in the product multiset.
\end{theorem}

\noindent For primes $q \geq 5$ where $|(\ZZ/q\ZZ)^\times| = q{-}1$ is
composite, non-faithful nontrivial characters exist, and the sole
remaining open hypothesis is
\defname{NonFaithfulCharacterEscape}---spectral gap non-summability
for characters with nontrivial kernel.  The faithful/non-faithful
decomposition is clean:

\begin{theorem}[\lean{EM/Advanced/VanishingNoiseVariant.lean}{nfce\_implies\_ufdStrong}]
\label{thm:nfce-ufdstrong}
$\mathrm{NonFaithfulCharacterEscape}(q)
\;\Longrightarrow\; \mathrm{UFDStrong}(q)$
for all primes $q \geq 3$.
\end{theorem}

\noindent The faithful half is free; only the non-faithful half
remains open.

\paragraph{NFCE infrastructure.}
The formalization develops a precise characterization of what NFCE
requires.  Define the \defname{factorRatio} at step~$n$ as
$\mathrm{liftToUnit}(\minFac) \cdot
\mathrm{liftToUnit}(\operatorname{secondMinFac})^{-1}$
in~$(\ZZ/q\ZZ)^\times$.
At non-fallback steps (raw card $\geq 2$), the spectral gap is zero
if and only if this ratio lies in~$\ker(\chi)$
(\lean{EM/Advanced/VanishingNoiseVariant.lean}{gap\_zero\_iff\_ratio\_in\_ker}).
NFCE thus reduces to: for each non-faithful nontrivial~$\chi$, the
factor ratio does not eventually stay in~$\ker(\chi)$.

If NFCE fails for a specific~$\chi$, the gaps are summable,
which forces \defname{KernelConfinement}: the ratio is eventually
confined to~$\ker(\chi)$
(\lean{EM/Advanced/VanishingNoiseVariant.lean}{summable\_implies\_ratio\_confined}).
Since $\ker(\chi)$ has index $\geq 2$ for nontrivial~$\chi$
(\lean{EM/Advanced/VanishingNoiseVariant.lean}{ker\_index\_ge\_two}),
kernel confinement is a strong arithmetic constraint:
$\minFac / \operatorname{secondMinFac} \bmod q$ eventually lies in a
proper subgroup of the unit group.

\paragraph{NFCE(5): reduction to a single quadratic escape.}
For $q = 5$, the unit group $(\ZZ/5\ZZ)^\times$ has order~$4$
(cyclic, $\cong \ZZ/4\ZZ$).  Its characters partition as: the trivial
character, two faithful characters (order~$4$, handled unconditionally
by FaithfulCharacterEscape), and one non-faithful nontrivial character
(order~$2$, the unique quadratic character).  Thus NFCE(5) reduces to a
single escape condition: the factor ratio
$\minFac / \operatorname{secondMinFac} \bmod 5$ must leave the unique
index-2 subgroup of $(\ZZ/5\ZZ)^\times$ infinitely often.
The formalization proves that this
\defname{RatioKernelEscape} condition suffices for the full chain
\lean{EM/Advanced/VanishingNoiseVariant.lean}{variant\_mc\_five\_of\_ratio\_escape}
{\color{leanink}$\checkmark$}:
RatioKernelEscape $\Rightarrow$ NFCE(5) $\Rightarrow$ UFDStrong(5)
$\Rightarrow$ StochasticTwoPointMC(5).

\paragraph{Intersection kernel dichotomy.}
When NFCE fails for \emph{all} non-faithful nontrivial characters
simultaneously (\defname{TotalNFCEFailure}), the factor ratio is
eventually confined to every non-faithful kernel (from the
summability-to-confinement lemma of Part~25).
The intersection of these kernels is constrained by a group-theoretic
fact: in a finite group whose order has $\geq 2$ distinct prime
factors, there exist subgroups of coprime orders whose intersection is
trivial
(\lean{EM/Advanced/VanishingNoiseVariant.lean}{coprime\_order\_zpowers\_inf\_bot}
{\color{leanink}$\checkmark$}).  Combined with Cauchy's theorem, this
yields \defname{NonFaithfulCharSeparation}~(NFCS): for every
$g \neq 1$, some non-faithful nontrivial character has $\chi(g) \neq 1$.
NFCS is
\textbf{proved} for all groups whose order has $\geq 2$ distinct prime
factors
(\lean{EM/Advanced/VanishingNoiseVariant.lean}{nonFaithfulCharSeparation\_of\_two\_prime\_factors}
{\color{leanink}$\checkmark$}), which for $(\ZZ/q\ZZ)^\times$ covers
all primes except $q = 2, 3, 5$ and Fermat primes (where $q - 1$ is a
power of~$2$).
Under NFCS, TotalNFCEFailure forces
$\minFac / \operatorname{secondMinFac} \equiv 1 \pmod{q}$ eventually
(\lean{EM/Advanced/VanishingNoiseVariant.lean}{total\_confinement\_implies\_ratio\_one}
{\color{leanink}$\checkmark$}).
The resulting \textbf{dichotomy}: for each qualifying~$q$, either
NFCE($q$) holds (giving UFDStrong), or the two smallest prime factors
become congruent mod~$q$ from some point on---a strong arithmetic
rigidity.

\paragraph{Comparison of open hypotheses.}
Standard MC requires DSL (population-to-orbit transfer for $\minFac$
residues).  The $\{\minFac, \operatorname{secondMinFac}\}$ variant
requires UFDStrong (spectral diversity of two-point factor sets), and
for $q = 3$ this is now \textbf{unconditional}; for $q \geq 5$ the
residual hypothesis is NonFaithfulCharacterEscape.
Neither DSL nor UFDStrong implies the other: DSL is about the
\emph{selected} factor's equidistribution; UFDStrong is about the
\emph{pair}'s spectral gap.
But both are instances of the same meta-question: does the deterministic
EM construction produce enough diversity in its factorizations to
prevent persistent alignment with any fixed character?

\paragraph{Ensemble two-point: population-level reduction.}
The ensemble framework of Section~\ref{sec:ensemble} can also be
applied to the two-point variant.  For generalized EM starting from
squarefree~$n$, the \defname{genFactorRatio} at step~$k$ is the ratio
$\minFac / \operatorname{secondMinFac} \bmod q$ in
$(\ZZ/q\ZZ)^\times$.  For a nontrivial character~$\chi$, a step
``escapes'' if this ratio lies outside~$\ker(\chi)$---producing a
spectral gap.  The formalization
(\code{Ensemble/TwoPointEnsemble.lean}, 566~lines, zero \code{sorry})
develops the full first-moment chain:
\[
  \text{PRED}
  \;\xrightarrow[\text{Fubini}]{\checkmark}\;
  \text{linear first moment}
  \;\xrightarrow[\text{partition}]{\checkmark}\;
  \text{positive density of high escapers.}
\]
Here \textbf{PopulationRatioEscapeDensity}~(PRED) asserts that for
each nontrivial~$\chi$ and each step $k \geq 1$, a positive
density~$\delta > 0$ of squarefree starting points have factor ratio
outside~$\ker(\chi)$.  Under PRED, the ensemble average of cumulative
escape counts grows linearly
(\lean{EM/Ensemble/TwoPointEnsemble.lean}{escape\_first\_moment\_linear}
{\color{leanink}$\checkmark$}), and the partition argument gives
\lean{EM/Ensemble/TwoPointEnsemble.lean}{positive\_density\_high\_escape}
{\color{leanink}$\checkmark$}: for each threshold~$M$, a positive
density of starting points achieve $\geq M$ cumulative escapes
within~$K$ steps.  PRED itself is open; it should follow from
MFRE (MinFacResidueEquidist) via CRT independence of the two
smallest factors.

\paragraph{Self-consistent random walk: the Fourier bridge.}
The stochastic two-point variant above fixes the factor sets
$\{p_1(n), p_2(n)\}$ from the \emph{actual} EM walk and randomizes
the selection.  A stronger variant makes the construction
\emph{self-consistent}: at each step, the accumulator determines the
available factors, and the fair-coin choice feeds back into the next
accumulator.  Formally, define
$\mathrm{epsWalkProdFrom}(\mathrm{acc}, \sigma, 0) = \mathrm{acc}$,
$\mathrm{epsWalkProdFrom}(\mathrm{acc}, \sigma, n{+}1)
= P_n \cdot (\text{if } \sigma(n) \text{ then } \minFac(P_n{+}1)
\text{ else } \operatorname{secondMinFac}(P_n{+}1))$
where $P_n = \mathrm{epsWalkProdFrom}(\mathrm{acc}, \sigma, n)$.
The decision sequence $\sigma : \NN \to \mathrm{Bool}$ replaces the
deterministic $\minFac$ rule with a random binary choice.

The central result is the \textbf{Fourier bridge identity}
(\lean{EM/Advanced/RandomTwoPointMC.lean}{fourier\_bridge\_identity\_proved}
{\color{leanink}$\checkmark$}): for any accumulator $\mathrm{acc} \geq 2$,
\[
  \mathrm{treeCharSum}(q, \chi, N, \mathrm{acc}, \tfrac{1}{2})
  \;=\;
  \frac{1}{2^N}
  \sum_{\sigma \in \{0,1\}^N}
  \prod_{k < N} \mathrm{chiAt}(q, \chi,
    \mathrm{epsWalkFactorFrom}(\mathrm{acc}, \sigma, k)).
\]
The proof is by induction on~$N$: at each step, the $2^N$ paths split
by first bit into two halves corresponding to the left and right
subtrees of the tree recursion; shift lemmas
(\code{epsWalkProdFrom\_shift}, \code{epsWalkFactorFrom\_shift})
connect the concatenated walk to the shifted walk at the new
accumulator; the $\tfrac{1}{2}$ weights combine to match the tree
recurrence.  Specializing to $\mathrm{acc} = 2$ gives the
unconditional identity
$\mathrm{fairTreeCharSum} = \mathrm{pathCharSum}$
(\lean{EM/Advanced/RandomTwoPointMC.lean}{fairTreeCharSum\_eq\_pathCharSum}
{\color{leanink}$\checkmark$}).

\defname{RandomTwoPointMC}
(\lean{EM/Advanced/RandomTwoPointMC.lean}{RandomTwoPointMC})
asks: for $q \geq 3$ and every unit $a \in (\ZZ/q\ZZ)^\times$, does
some binary decision sequence of finite length produce a walk from
$\mathrm{acc} = 2$ landing in residue class~$a$?  This is the
self-consistent analog of StochasticTwoPointMC.
The Fourier bridge connects the tree to the path character sum.
To connect path character sums to reached residue classes, one needs
\defname{ChiAtMultiplicativity}: that $\mathrm{chiAt}(\mathrm{acc})
\cdot \prod_k \mathrm{chiAt}(\text{factor}_k)
= \mathrm{chiAt}(\text{endpoint})$, provided $q$ does not divide the
endpoint.  This is \textbf{proved}
(\lean{EM/Advanced/RandomTwoPointMC.lean}{chiAt\_multiplicativity\_proved}
{\color{leanink}$\checkmark$})
by induction on~$N$ with backward non-divisibility propagation:
if $q \nmid P_N$, then $q \nmid P_k$ for all $k \leq N$
(\lean{EM/Advanced/RandomTwoPointMC.lean}{not\_dvd\_epsWalkProdFrom\_of\_le}
{\color{leanink}$\checkmark$}),
so every factor is coprime to~$q$ and $\mathrm{chiAt}$ is
multiplicative.

A full Fourier counting infrastructure is built on top: character
orthogonality on $(\ZZ/q\ZZ)^\times$, death/survival partition
(\lean{EM/Advanced/RandomTwoPointMC.lean}{survival\_plus\_death}
{\color{leanink}$\checkmark$}:
$\mathrm{survivalCount} + \mathrm{deathCount} = 2^N$),
the Fourier indicator formula
$|G| \cdot \mathrm{pathCount}(a)
= \sum_\chi \overline{\chi(a)} \cdot U(\chi, N)$,
and norm bounds for the unit endpoint character sum.
The chain is:
\[
  \mathrm{TCA}
  \;\xRightarrow{\checkmark}\;
  \text{vanishing pathCharSum}
  \;\xRightarrow{\text{ChiAtMult}\ \checkmark}\;
  \text{Fourier counting}
  \;\xRightarrow{?}\;
  \text{RandomTwoPointMC}.
\]
The remaining gap is \defname{PathSurvival}
(\lean{EM/Advanced/RandomTwoPointMC.lean}{PathSurvival}):
the ratio $\mathrm{survivalCount}(N) / \mathrm{deathCount}(N)$
grows without bound.  Without positive survival at the same~$N$
where pathCharSum is small, the Fourier inversion gives no
information about unit classes.

For $q = 3$, the analysis reveals a structural obstruction.
Since $\mathrm{acc} = 2$ and $\minFac(3) = \operatorname{secondMinFac}(3) = 3$,
the walk immediately dies: $3 \mid P_1 = 6$ on \emph{all} paths,
and absorption is permanent.
Thus PathSurvival(3) is false---but
$\mathrm{TreeContractionAtHalf}(3)$ is \emph{also} false (the tree
character sum is identically~$1$ at every depth, because all
post-death factors are $\equiv 1 \pmod{3}$), making
$\mathrm{TreeContractionImpliesRandomMC}(3)$ vacuously true.

For $q = 3$, the \textbf{zero-mean property} gives maximal per-step
contraction
(\lean{EM/Advanced/RandomTwoPointMC.lean}{mean\_char\_zero\_at\_diverse\_step\_three}
{\color{leanink}$\checkmark$}):
$(1/2)\chi(a) + (1/2)\chi(b) = 0$ whenever $\chi(a) \neq \chi(b)$
for a nontrivial character on $(\ZZ/3\ZZ)^\times$ (order~$2$, so
$\chi$ maps to $\{1, -1\}$).  Although this property is proved,
the immediate death at $q = 3$ means the zero-mean steps are
never reached from the standard starting point.

\paragraph{TCA $+$ PathSurvival $\Rightarrow$ RandomTwoPointMC.}
The conditional implication is now \textbf{proved}
(\lean{EM/Advanced/RandomTwoPointMC.lean}{tca\_path\_survival\_implies\_random\_mc\_proved}
{\color{leanink}$\checkmark$}).
The proof~(${\sim}400$~lines) builds the full Fourier counting
infrastructure on $(\ZZ/q\ZZ)^\times$:
\begin{enumerate}[nosep]
\item \textbf{Path sum decomposition}
  (\lean{EM/Advanced/RandomTwoPointMC.lean}{path\_sum\_decomp}
  {\color{leanink}$\checkmark$}):
  the unscaled path sum
  $2^N \cdot \mathrm{pathCharSum}$ splits as
  $\mathrm{chiAt}(2)^{-1}\cdot U(\chi,N) + E(\chi,N)$,
  where $U$ sums over surviving paths via ChiAtMultiplicativity and
  $E$ sums over death paths.

\item \textbf{Death error bound}
  (\lean{EM/Advanced/RandomTwoPointMC.lean}{death\_error\_norm\_le}
  {\color{leanink}$\checkmark$}):
  $\|E(\chi,N)\| \leq \mathrm{deathCount}(N)$
  (each death term has norm $\leq 1$).

\item \textbf{Unit endpoint bound}
  (\lean{EM/Advanced/RandomTwoPointMC.lean}{uecsi\_norm\_le}
  {\color{leanink}$\checkmark$}):
  $\|U(\chi,N)\| \leq 2^N\|\mathrm{pathCharSum}\| + \mathrm{deathCount}$.

\item \textbf{Fourier counting identity}
  (\lean{EM/Advanced/RandomTwoPointMC.lean}{fourier\_counting\_identity}
  {\color{leanink}$\checkmark$}):
  $\sum_\chi \overline{\chi(a)} \cdot U(\chi,N)
  = |G| \cdot \mathrm{pathCount}(a,N)$.

\item \textbf{Fourier lower bound}
  (\lean{EM/Advanced/RandomTwoPointMC.lean}{fourier\_lower\_bound}
  {\color{leanink}$\checkmark$}):
  $(q{-}1)\cdot\mathrm{pathCount}(a) \geq
  \mathrm{survivalCount} - (q{-}2)\cdot B$
  where $B$ is the uniform bound on nontrivial $\|U(\chi)\|$.
\end{enumerate}
Combining: TCA gives $\|\mathrm{pathCharSum}\| < 1/(2q)$ for
$N \geq N_0$; PathSurvival gives
$\mathrm{survivalCount} > 2q^2 \cdot \mathrm{deathCount}$ for
$N \geq N_1$.  At $N = \max(N_0,N_1)$, the Fourier lower bound yields
$\mathrm{pathCount}(a) > 0$ for every unit~$a$, via the algebraic
identity $2q^3 + 2q^2 + 3q + 2 > 0$ (true for all $q \geq 5$).
For $q = 3$, TCA is vacuously false, so the implication holds trivially.
The self-consistent RandomTwoPointMC thus reduces to exactly two open
hypotheses: TreeContractionAtHalf and PathSurvival.

\paragraph{Ensemble $\varepsilon$-walk.}
The formalization
(\code{Advanced/EpsilonWalk.lean}, 458~lines, zero \code{sorry})
develops an ensemble version: instead of fixing $\mathrm{acc} = 2$,
parametrize the $\varepsilon$-walk by squarefree starting point~$n$.
The \defname{SigmaCRTPropagationStep} hypothesis (open) asserts that
CRT equidistribution propagates across steps for $\sigma$-walks,
mirroring CRTPropagationStep for deterministic walks.
Under SRE $+$ SigmaCRTPropagationStep, the death density (fraction of
squarefree~$n$ with $q \mid \mathrm{epsWalkProdFrom}(n,\sigma,k)+1$)
converges to~$1/(q{-}1)$ for every fixed~$k$ and~$\sigma$
(\lean{EM/Advanced/EpsilonWalk.lean}{sigma\_death\_density\_tendsto}
{\color{leanink}$\checkmark$}).

\subsection{Ensemble MC: Almost Every Starting Point}
\label{sec:ensemble-variant}

The variants above change the \emph{selection rule}; here we change the
\emph{starting point}.  The generalized EM construction
(Section~\ref{sec:ensemble}) parametrizes the sequence by its initial
value~$n$: $\Prod{n}(0) = n$,
$\seq{n}(k) = \minFac(\Prod{n}(k)+1)$,
$\Prod{n}(k{+}1) = \Prod{n}(k)\cdot\seq{n}(k)$.
The standard orbit is $n = 2$.

\begin{definition}[\lean{EM/Ensemble/WeylChain.lean}{GenMullinConjecture}]
$\mathrm{GenMC}(n)$: every prime appears in the generalized EM
sequence from~$n$, i.e.\ $\forall q$ prime, $\exists k$,
$\seq{n}(k) = q$.
\end{definition}

The reduction
\lean{EM/Ensemble/WeylChain.lean}{gen\_hitting\_implies\_gen\_mc\_proved}
{\color{leanink}$\checkmark$} proves that cofinal walk hitting
implies~$\mathrm{GenMC}(n)$ for every squarefree~$n$, mirroring the
standard inductive bootstrap.  And
\lean{EM/Ensemble/WeylChain.lean}{gen\_mc\_two\_implies\_mc}
{\color{leanink}$\checkmark$} shows $\mathrm{GenMC}(2) \Rightarrow
\mathrm{MC}$, so the generalized conjecture specializes to the
standard one.

\paragraph{Bag arithmetic.}
At each step~$k$, the Euclid number $\Prod{n}(k)+1$ has a ``bag'' of
prime factors; the selected prime $\seq{n}(k)$ is the minimum of this
bag.  The formalization
(\code{Ensemble/BagArithmetic.lean}, 224~lines, zero \code{sorry})
develops the bag-level quantities:
\defname{genEuclidOmega}~$\omega(n,k) = |\text{primeFactors}(\Prod{n}(k)+1)|$
(at least~1,
\lean{EM/Ensemble/BagArithmetic.lean}{genEuclidOmega\_pos}
{\color{leanink}$\checkmark$}),
\defname{genBagDiversity} (number of distinct residue classes mod~$q$
among the factors, $\leq \omega$ and $\leq q$),
\defname{genFactorsInClass}~$F_a$ (factors in class~$a$, with the
partition identity
$\sum_a |F_a| = \omega$,
\lean{EM/Ensemble/BagArithmetic.lean}{genFactorsInClass\_card\_sum}
{\color{leanink}$\checkmark$}),
and \defname{genEuclidCofactor} (the quotient
$(\Prod{n}(k)+1)/\seq{n}(k)$, whose prime factors are a subset of
the bag,
\lean{EM/Ensemble/BagArithmetic.lean}{cofactor\_primeFactors\_subset}
{\color{leanink}$\checkmark$}).
For $k \geq 1$, all bag elements are $\geq 3$
(\lean{EM/Ensemble/BagArithmetic.lean}{primeFactors\_ge\_three\_of\_pos}
{\color{leanink}$\checkmark$}).

\paragraph{Density variants.}
Instead of proving $\mathrm{GenMC}(n)$ for the single orbit $n = 2$,
one can ask how many starting points satisfy it.  Two density levels
are formalized:

\begin{itemize}
\item \textbf{PositiveDensityRSD}
  (\lean{EM/Population/ReciprocalSum.lean}{PositiveDensityRSD}):
  a positive density~$\delta > 0$ of squarefree~$n$ have divergent
  reciprocal sums $\sum_k 1/\seq{n}(k) = \infty$.

\item \textbf{AlmostAllSquarefreeRSD}
  (\lean{EM/Population/ReciprocalSum.lean}{AlmostAllSquarefreeRSD}):
  density~$1$ of squarefree~$n$ have divergent reciprocal sums.
\end{itemize}

\noindent
The reciprocal sum is a natural proxy: if $\sum 1/\seq{n}(k)$
diverges, then the primes appearing in the sequence from~$n$ cannot be
confined to a thin set; although divergence alone does not formally
imply~$\mathrm{GenMC}(n)$, it rules out the worst obstruction (a
bounded sequence of captured primes).

\paragraph{The one-hypothesis route: LMG.}
The proved chain
\lean{EM/Population/ReciprocalSum.lean}{lmg\_implies\_positive\_density\_rsd}
{\color{leanink}$\checkmark$} shows:
\[
  \mathrm{LMG} \;\Longrightarrow\; \mathrm{PositiveDensityRSD},
\]
where \textbf{LinearMeanGrowth}~(LMG) asks that
$\mathbb{E}_n[S_K(n)] \geq \kappa K$ for some $\kappa > 0$ and all
$K$ large enough (here $S_K(n) = \sum_{k<K} 1/\seq{n}(k)$ is the
partial reciprocal sum, averaged over squarefree~$n$).
LMG is a \emph{population-level} statement: it involves no specific
orbit, bypassing the orbit-specificity barrier that obstructs DSL\@.

With the additional hypothesis PairwiseStepDecorrelation~(PSD)---that
$\operatorname{Cov}[1/\seq{n}(j),\, 1/\seq{n}(k)] \to 0$ for
$j \neq k$---the stronger conclusion follows:
$\mathrm{PSD} + \mathrm{LMG} \Rightarrow \mathrm{AlmostAllSquarefreeRSD}$
(\lean{EM/Population/ReciprocalSum.lean}{dead\_end\_125\_pairwise\_revival}
{\color{leanink}$\checkmark$}).
The bridges are all machine-verified:
$\mathrm{IVB}(1/4)$ {\color{leanink}$\checkmark$},
$\mathrm{PSD} + \mathrm{IVB} \Rightarrow$ variance bound
{\color{leanink}$\checkmark$},
Chebyshev concentration {\color{leanink}$\checkmark$}.

\paragraph{CRT propagation and the first moment.}
LMG follows from
\textbf{FirstMomentStep}~(FMS):
$\mathbb{E}_n[1/\seq{n}(k)] \to \kappa$ as $X \to \infty$ for each~$k$
(\lean{EM/Ensemble/FirstMoment.lean}{FirstMomentStep}).
FMS in turn follows from \textbf{CRTPropagationStep}
(\lean{EM/Ensemble/CRT.lean}{CRTPropagationStep}):
the equidistribution of $\Prod{n}(k) \bmod q$ (among coprime residues)
propagates from step~$k$ to step~$k{+}1$.
The chain is:
\[
  \mathrm{SRE}
  \;+\; \mathrm{CRTPropagationStep}
  \;\xRightarrow{\text{\color{leanink}$\checkmark$}}\;
  \mathrm{AccumEquidist}
  \;\Longrightarrow\;
  \mathrm{FMS}(\kappa)
  \;\Longrightarrow\;
  \mathrm{LMG}
  \;\xRightarrow{\text{\color{leanink}$\checkmark$}}\;
  \mathrm{PositiveDensityRSD}.
\]
Here SRE (squarefree residue equidistribution) is the standard base
case, and AccumEquidist~$\Rightarrow$~FMS uses population
equidistribution of $\minFac$ via the Alladi--Buchstab weights.

CRTPropagationStep now reduces to a cleaner open hypothesis:
\textbf{EnsembleTransitionApprox}~(ETA,
\lean{EM/Ensemble/BackwardDynamics.lean}{EnsembleTransitionApprox}),
which says that the empirical transition probability---the fraction of
squarefree~$n$ with $\Prod{n}(k) \equiv c$ whose multiplier satisfies
$\seq{n}(k) \equiv b$---converges to $1/(q{-}1)$ for all nonzero
$b, c$.  The proved reduction
\lean{EM/Ensemble/BackwardDynamics.lean}{eta\_implies\_crt\_propagation}
{\color{leanink}$\checkmark$} shows
$\mathrm{ETA} \Rightarrow \mathrm{CRTPropagationStep}$
via the backward density decomposition: the density at step $k{+}1$ is
the matrix product of the density at step~$k$ with the transition
matrix, and when both converge to~$1/(q{-}1)$, so does the product.
The full chain
\lean{EM/Ensemble/BackwardDynamics.lean}{eta\_sre\_implies\_prsd}
{\color{leanink}$\checkmark$} composes this with the proved bridges to
give $\mathrm{ETA} + \mathrm{SRE} \Rightarrow
\mathrm{PositiveDensityRSD}$.

The structural basis for ETA is CRT multiplier invariance
(\lean{EM/Group/CRT.lean}{crt\_multiplier\_invariance}
{\color{leanink}$\checkmark$}): $\minFac(\Prod{n}(k)+1)$ depends on
the full integer, not just its residue mod~$q$, so the factoring
outcome is ``$q$-blind.''

\paragraph{Backward dynamics approach (Tao--Collatz analogy).}
The backward dynamics infrastructure is formalized in
\code{Ensemble/BackwardDynamics.lean} (494~lines).
Joint counts, transition probabilities, their basic probability axioms
(non-negativity, partition identity, sum-to-one), and the backward
density decomposition are all proved.
The approach follows Tao's backward
dynamics for the Collatz conjecture~\cite{Tao2019Collatz}: for a
target residue $a \in (\ZZ/q\ZZ)^\times$ at step~$k{+}1$,
\[
  F_{k+1}(a) = \sum_{c \in (\ZZ/q\ZZ)^\times} F_k(c) \cdot R(c, a, k, X),
\]
where $R(c, a, k, X) = 1/(q{-}1) + O(1/\log X)$ by CRT-blindness.
When $F_k$ is already uniform, so is $F_{k+1}$, giving the induction
step.  The key question is whether the error $O(1/\log X)$ accumulates
across~$k$ (contraction vs.\ expansion of the deviation vector);
super-exponential growth of $\Prod{n}(k)$ suggests contraction.

\paragraph{Comparison with DSL.}
The DSL route (Section~\ref{sec:ensemble}) proves MC for the single
orbit $n = 2$, contingent on transferring population statistics to one
deterministic trajectory.
The ensemble route here proves a density result for \emph{almost all}
starting points, contingent on CRTPropagationStep, which is a
population-level hypothesis with no orbit-specificity.
Neither subsumes the other: DSL would give MC directly; CRTPropagationStep
gives PositiveDensityRSD (which does not formally imply MC for $n = 2$).
However, the ensemble route faces a structurally easier barrier, since
its open hypothesis is a statement about averages, not about one
specific orbit.

\subsection{Weak Mullin and Analytic Weakenings}
\label{sec:weak-mullin}

The variants above change the \emph{selection rule}; this subsection
changes the \emph{conclusion}.  Instead of proving that every prime
appears ($M = \emptyset$), one can ask that the set of missing
primes is small, or that the EM reciprocal sum diverges.

\paragraph{Missing primes and Weak Mullin.}
Define the set of \emph{missing primes}
\[
M = \{p \in \mathbb{P} : p \neq \seq{2}(k) \text{ for all } k\}.
\]
MC is the assertion $M = \emptyset$.  A natural weakening asks only that
$M$ is \emph{small} in an analytic sense:

\begin{property}[\defname{WeakMullin} (\abbr{WM})]
\lean{EM/Population/WeakMullin.lean}{WeakMullin}
The reciprocal sum over missing primes converges:
$\sum_{q \in M} 1/q < \infty$.
\end{property}

This is strictly weaker than MC (which gives $M = \emptyset$ and hence
a trivially convergent empty sum) but strong enough to have a nontrivial
consequence:

\begin{property}[\defname{ReciprocalDivergence} (\abbr{RD})]
\lean{EM/Population/WeakMullin.lean}{ReciprocalDivergence}
The reciprocal sum over EM primes diverges:
$\sum_{k=0}^{\infty} 1/\seq{2}(k) = \infty$.
\end{property}

\begin{theorem}[\lean{EM/Population/WeakMullin.lean}{mc\_implies\_wm}]
$\MC \Longrightarrow \mathrm{WM}$.
\end{theorem}

\begin{theorem}[\lean{EM/Population/WeakMullin.lean}{wm\_implies\_rd}]
\label{thm:wm-rd}
$\mathrm{WM} \Longrightarrow \mathrm{RD}$.
\end{theorem}

\noindent
The proof partitions $\mathbb{P} = A \cup M$ (appeared vs.\ missing),
notes that both $\Sigma_A$ and $\Sigma_M$ converge under WM $+$ $\neg$RD,
contradicting the divergence of $\sum_{p} 1/p$.

\paragraph{Product divergence.}
The negation of WM has a quantitative consequence: $\neg\mathrm{WM}$ gives
$\sum_{q \in M} 1/(q{-}1) = \infty$, so the ``safe fraction'' product
$\prod_{q \in F}(q{-}2)/(q{-}1) \leq e^{-\sum 1/(q-1)} \to 0$ and its
reciprocal $\prod (q{-}1)/(q{-}2)$ exceeds any bound
(\lean{EM/Population/WeakMullin.lean}{not\_wm\_product\_diverges}).
The product $\prod (q{-}1)/(q{-}2)$ over missing primes measures the
``confinement cost'': each missing prime~$q$ confines the walk to $q-2$
of the $q-1$ units in $(\ZZ/q\ZZ)^\times$ (avoiding~$-1$).  If WM
fails, this confinement cost is unbounded.

\paragraph{Confined energy.}
For a missing prime $q \geq 3$, the walk avoids $-1$ entirely, so the
$N$ visits are distributed among only $q-2$ residues.  A pigeonhole
argument gives a confined-energy lower bound
(\lean{EM/Population/WeakMullin.lean}{per\_prime\_confined\_energy\_lb}):
$N^2/(q{-}2) \leq (q{-}1)\sum_a V_N(a)^2$.
The excess energy ratio $(q{-}1)/(q{-}2)$ compared to the
equidistributed baseline $N^2/(q{-}1)$ is bounded away from~$1$.
\defname{JointSVE}
(\lean{EM/Population/WeakMullin.lean}{JointSVE})---joint subquadratic
visit energy over any finite set of missing primes---implies MMCSB
(\lean{EM/Population/WeakMullin.lean}{jointSVE\_implies\_mmcsb}),
connecting the ``missing primes'' picture to the spectral route.

\paragraph{The shifted squarefree population.}
The EM accumulator $\Prod{2}(n)$ is squarefree
(\lean{EM/Population/WeakErgodicity.lean}{prod\_squarefree}),
so every Euclid number $\Prod{2}(n)+1$ belongs to the shifted
squarefree population
$\mathcal{S} = \{m \geq 2 : m{-}1 \text{ squarefree}\}$
(density $6/\pi^2$).  This decomposes the orbit equidistribution
problem into:
\begin{itemize}[nosep]
\item \textbf{Population Equidistribution (PE):} among elements of
  $\mathcal{S}$, the smallest prime factor is equidistributed modulo
  any prime~$q$ (provable by Selberg sieve $+$ Dirichlet).
\item \textbf{Population Transfer (PT):} the EM trajectory samples
  $\mathcal{S}$ without introducing mod-$q$ bias.
  This is the open content.
\end{itemize}
The chain PE $+$ PT $+$ EMDImpliesCME $\Longrightarrow$ MC is proved
(\lean{EM/Population/WeakErgodicity.lean}{pe\_transfer\_cme\_implies\_mc}).

\paragraph{Pairwise decorrelation and positive density.}
Dead End~\#125 proved that pairwise cancellation does not imply $k$-wise,
but \emph{variance bounds only use $k=2$}.  The formalization proves the
revival chain
(\lean{EM/Population/ReciprocalSum.lean}{dead\_end\_125\_pairwise\_revival}):
PairwiseStepDecorrelation $+$ LinearMeanGrowth $+$
ChebyshevConcentration $\Longrightarrow$ AlmostAllSquarefreeRSD\@.
A further simplification drops PSD entirely:

\begin{theorem}[\lean{EM/Population/ReciprocalSum.lean}{lmg\_implies\_positive\_density\_rsd}]
\label{thm:positive-density-rsd}
$\mathrm{LinearMeanGrowth} \;\Longrightarrow\; \mathrm{PositiveDensityRSD}$.
\end{theorem}

\noindent This is a 1-hypothesis route: LMG is strictly weaker than
CME---it requires only that the ensemble mean grows linearly, not
conditional equidistribution.

\subsection{The Function Field Analog}
\label{sec:function-field}

The Euclid--Mullin construction makes sense over any unique
factorization domain.  Replacing $\ZZ$ by $\FF_p[t]$ and primes by
monic irreducible polynomials gives a sequence
\[
  f_0 = t, \qquad
  f_{n+1} = \text{min-degree monic irreducible factor of }
  F(n) + 1,
\]
where $F(n) = f_0 \cdots f_n$ is the running product.  We call this the
\textbf{FF-EM sequence} and state the function field analog of MC\@:

\begin{conjecture}[\defname{FF-MC}]
\lean{EM/FunctionField/Analog.lean}{FFMullinConjecture}
Every monic irreducible polynomial in $\FF_p[t]$ eventually appears in
the FF-EM sequence.
\end{conjecture}

The function field setting is instructive because the Weil bound---an
\emph{open hypothesis} over $\ZZ$ (Wiener--Ikehara + Alladi)---is a
\emph{theorem} over $\FF_p[t]$.  This eliminates the analytic number
theory entirely: population equidistribution is free.  The question
becomes: what remains after removing the ANT difficulty?

\paragraph{What the Weil bound gives.}
Let $Q \in \FF_p[t]$ be a monic irreducible of degree~$d$, and
$\chi : (\FF_p[t]/(Q))^\times \to \mathbb{C}^\times$ a non-trivial
multiplicative character.  The Weil bound for function fields states:
\[
  \Bigl|\sum_{\substack{\alpha \in \FF_p[t] \\ \deg \alpha < d}}
  \chi(\alpha \bmod Q)\Bigr|
  \;\leq\; (d - 1)\sqrt{p}.
\]
This immediately gives \textbf{population equidistribution}
(\lean{EM/FunctionField/Analog.lean}{weil\_implies\_population\_equidist}):
among all monic irreducibles, the residues modulo~$Q$ are
equidistributed in $(\FF_p[t]/(Q))^\times$.  Over $\ZZ$, obtaining
the analogous statement requires the Weighted PNT in Arithmetic
Progressions (= Wiener--Ikehara), which is the sole unformalized ANT
input.  Over $\FF_p[t]$, it is unconditional.

\paragraph{What the Weil bound does not give.}
Population equidistribution controls character sums over
\emph{all} monic irreducibles---a generic sum.  FF-MC requires
controlling the factorization of the \emph{specific} polynomial
$F(n) + 1$, which is determined by the entire history of the
sequence.  The Weil bound says nothing about which irreducible
factor of a specific polynomial has the smallest degree.

This is the \textbf{orbit-specificity barrier} (Dead End~\#127,
\lean{EM/FunctionField/Analog.lean}{WeilDoesNotImplyFFDSL}),
structurally identical to the integer case: no known result transfers
population-level equidistribution to a specific deterministic orbit.
The function field analog thus confirms that the difficulty of MC is
\emph{not} the analytic number theory but the orbit transfer.

\paragraph{The bootstrap: competitor exhaustion.}
Over $\ZZ$, the bootstrap induction is clean: $\minFac$ selects a
unique prime, so if $q \mid \Prod{2}(n)+1$ and all smaller primes
have appeared, then $\seq{2}(n{+}1) = q$.  Over $\FF_p[t]$, multiple
irreducibles of the same degree can divide $F(n)+1$, and the
min-degree rule only guarantees that the selected factor has degree
$\leq \deg Q$.

The solution is \textbf{competitor exhaustion}
(\lean{EM/FunctionField/Bootstrap.lean}{ff\_element\_capture}):
there are finitely many monic irreducibles of each degree~$d$ over
$\FF_p$ (exactly $\frac{1}{d}\sum_{e \mid d} \mu(d/e)\, p^e$ of
them).  If $Q$ divides $F(n)+1$ cofinally, the hitting set
$\{n : Q \mid F(n)+1\}$ is infinite.  At each such~$n$, the
selected factor $f_{n+1}$ has $\deg f_{n+1} \leq \deg Q$ and is
injective into the finite set of competitors
(by \lean{EM/FunctionField/Bootstrap.lean}{ffSeq\_injective\_proved}).
Since an injection from an infinite set to a finite set is
impossible, eventually all competitors of degree $\leq \deg Q$
are consumed, and $Q$ must be selected.

\begin{theorem}[\lean{EM/FunctionField/Bootstrap.lean}{ff\_dh\_implies\_ffmc}]
\label{thm:ff-bootstrap}
$\mathrm{FFDynamicalHitting} \;\Longrightarrow\; \mathrm{FF\text{-}MC}$.
\end{theorem}

Here \defname{FFDynamicalHitting} asserts: for every monic irreducible
$Q$ not in the sequence, $Q \mid F(n)+1$ for cofinally many~$n$.
The proof is a strong induction on $\deg Q$: given FF-MC below
degree~$d$, cofinal hits by~$Q$ plus competitor exhaustion force~$Q$
to appear.

\paragraph{SubgroupEscape via Weil.}
Over $\ZZ$, SubgroupEscape (SE) is proved elementarily via
PrimeResidueEscape.  Over $\FF_p[t]$, the Weil bound gives a
quantitative version for \emph{large~$p$}.

The idea: suppose all degree-1 multipliers $\{(t - a) \bmod Q : a \in \FF_p\}$
lie in a proper subgroup $H \leq (\FF_p[t]/(Q))^\times$ of index~$\ell$.
Then for any character~$\chi$ that kills~$H$ (i.e.,
$\chi(h) = 1$ for $h \in H$), the character sum $\sum_{a \in \FF_p}
\chi((t-a) \bmod Q) = p$ since every term is~$1$.  But the Weil bound
gives $|\text{sum}| \leq (d-1)\sqrt{p}$, yielding $p \leq (d-1)\sqrt{p}$,
i.e., $p \leq (d-1)^2$.

\begin{theorem}[\lean{EM/FunctionField/SubgroupEscape.lean}{weil\_implies\_ff\_se}]
\label{thm:ff-se}
If $p \geq (d-1)^2 + 1$, then the degree-1 multipliers generate
the full unit group $(\FF_p[t]/(Q))^\times$ for every monic
irreducible~$Q$ of degree~$d$.
\end{theorem}

The threshold $(d-1)^2 + 1$ is sharp for this argument.
For $d = 1$, SE holds for all~$p$ (the degree-1 images
\emph{are} all nonzero elements of~$\FF_p$,
\lean{EM/FunctionField/SubgroupEscape.lean}{ff\_se\_degree\_one}).
For $d = 2$, the threshold is~$2$, so SE holds for all odd
primes.

\paragraph{The generation-to-coverage gap.}
SE proves that the multipliers \emph{generate} the full unit group.
Does this imply the walk \emph{visits}~$-1$?  No.

\begin{example}[\lean{EM/FunctionField/CyclicWalkCoverage.lean}{alternating\_walk\_misses\_two}]
\label{ex:z4}
In $\ZZ/4\ZZ$, the alternating steps $\{1, 3, 1, 3, \ldots\}$
generate the full group ($\gcd(1,3) = 1$), but the additive walk
$0 \to 1 \to 0 \to 1 \to \cdots$ visits only $\{0, 1\}$, permanently
missing $2$ and~$3$.
\end{example}

This is \textbf{Dead End~\#130}: SE does not imply DH.  Generation is
an algebraic property (the subgroup generated by the step set is
maximal); coverage is a dynamical property (the walk trajectory
actually reaches every element).  The two coincide for random walks
(Diaconis--Shahshahani) but diverge for deterministic ones.

\paragraph{The three-level hierarchy.}
The FF-MC reduction decomposes the orbit-specificity barrier into
three levels of difficulty
(\lean{EM/FunctionField/MultiplierCCSB.lean}{ff\_mc\_mult\_ccsb\_landscape}):

\begin{enumerate}[nosep]
\item \textbf{Population level} (solved).  Among \emph{all} monic
  irreducibles of degree~$e$, character values are equidistributed.
  Free from the Weil bound.

\item \textbf{Fiber level} (open).  Among irreducible factors of a
  \emph{specific} polynomial $f$ lying in a specific residue class
  mod~$Q$, character values cancel.  This requires orbit-specific
  information: Deligne's equidistribution theorem (Weil~II) applies
  to algebraic families, not to the sequential FF-EM construction.

\item \textbf{Selection level} (open).  Among the irreducible factors
  of~$f$, the one of \emph{smallest degree} has unbiased character
  value.  This is \defname{SelectionBiasNeutral}: does the greedy
  min-degree rule preserve character cancellation?  It is a purely
  combinatorial question, independent of algebraic geometry.
\end{enumerate}

The selection level is the true core of the difficulty.  Even if the
Weil bound gives population equidistribution and Deligne gives
fiber-level equidistribution, the minimum-degree selection could
systematically bias character values.

\paragraph{The cyclotomic obstruction.}
A natural strategy for fiber-level information is the monodromy
approach: if the Galois group of $F(n)+1$ is ``large'' (e.g.,
contains the alternating group~$A_d$), then Deligne's equidistribution
theorem~\cite{Deligne1980} constrains the Frobenius conjugacy class distribution
and hence the factorization statistics.

\begin{remark}[Dead End~\#129,
\lean{EM/FunctionField/Analog.lean}{ff\_cyclotomic\_dead\_end}]
\label{rem:fflm-false}
The \textbf{FF Large Monodromy} hypothesis (Galois group of
$F(n)+1$ eventually contains~$A_d$) is \emph{structurally false}.
The product of all monic irreducibles of degree $\leq d$ over~$\FF_p$
is $t^{p^d} - t$ (a standard identity).  Adding~$1$ gives
$t^{p^d} - t + 1$, whose splitting field is contained in
$\FF_{p^N}$ for some~$N$: the Galois group is \emph{abelian}.
Since the FF-EM sequence greedily consumes all irreducibles of
each degree, $F(n)$ eventually divides $t^{p^d} - t$, and the
Galois group of $F(n)+1$ is forced to be abelian.  Abelian
monodromy recovers only Dirichlet-type character sums---i.e.,
population equidistribution, which is already free from Weil.
\end{remark}

This rules out the monodromy route to FF-MC entirely.  The Galois
structure of the greedy Euclid construction is fundamentally different
from that of a ``random'' polynomial of the same degree.

\paragraph{Summary.}
The function field analog isolates the difficulty of Mullin's
Conjecture with surgical precision.  Removing the analytic number
theory (by working over $\FF_p[t]$ where Weil replaces
Wiener--Ikehara) reveals that the entire difficulty lies in the
orbit-specificity transfer: from population statistics to a specific
deterministic orbit.  The three dead ends (\#127: Weil does not
imply FF-DSL; \#129: large monodromy is false; \#130: SE does not
imply DH) collectively show that the function field setting, despite
its algebraic advantages, faces the same fundamental barrier as the
integer case.

The formalization across five files (2{,}302~lines, zero
\texttt{sorry}) makes the orbit-specificity barrier visible as a
mathematical theorem, not merely a heuristic observation:
\emph{every} algebraic tool that is available over function fields
has been applied, and the residual is precisely the same open
hypothesis (DSL/CME) that blocks MC over~$\ZZ$.
